In the previous section, we have shown how to automatically insert
$\kw{reset}$ and $\kw{reuse}$ instructions that minimize the number of
memory allocations at execution time. We now discuss techniques we have been using
for taking advantage of this transformation when writing functional
code. Two fundamental questions when using this optimization are: Does
a $\kw{reuse}$ instruction now guard my constructor applications?
Given a $\kw{let}\ y = \kw{reset}\ x$ instruction, how often is $x$
not shared at runtime? We address the first question using a simple
static analyzer that when invoked by a developer, checks whether
\kw{reuse} instructions are guarding constructor applications in a
particular function.  This kind of analyzer is straightforward to
implement in Lean since our IR is a Lean inductive datatype. This kind
of analyzer is in the same spirit of the \emph{inspection-testing}
package available for GHC~\cite{breitner2018promise}.
We cope with the second question using
runtime instrumentation.  For each $\kw{let}\ y = \kw{reset}\ x$
instruction, we can optionally emit two counters that track how often
$x$ is shared or not. We have found these two simple techniques quite
useful when optimizing our own code. Here, we report one instance
that produced a significant performance improvement.
%% Here, we report two instances that produced significant performance improvements:
%% \emph{red-black trees} and \emph{monad transformers}.

\subsection*{Red-black trees}
Red-black trees are implemented in the Lean standard
library and are often used to write proof automation.
For the purposes of this section, it is sufficient to have
an abstract description of this kind of tree, and one of the re-balancing functions used by the insertion
function.
\begin{lstlisting}
Color  = R | B
Tree a = E | T Color (Tree a) a (Tree a)
balance_1 v t (T _ (T R l x r_1) y r_2)  = T R (T B l x r_1) y (T B r_2 v t)
balance_1 v t (T _ l_1 y  (T R l_2 x r)) = T R (T B l_1 y l_2) x (T B r v t)
balance_1 v t (T _ l y r)              = T B (T R l y r) v t
insert (T B a y b) x = balance_1 y b (insert a x) if x < y and a is red
...
\end{lstlisting}
Note that the first two $\textit{balance}_1$ equations create three $T$ constructor values, but the patterns on the
left-hand side use only two $T$ constructors. Thus, the generated IR for $\textit{balance}_1$ contains $T$ constructor
applications that are not guarded by $\textit{reuse}$, and this fact can be detected at compilation time.
Note that even if the result of $(\textit{insert}\ a\ x)$ contains only nonshared values, we still have to allocate one
constructor value. We can avoid this unnecessary memory allocation by inlining $\textit{balance}_1$.
After inlining, the input value $(T\ B\ a\ y\ b)$ is reused in the $\textit{balance}_1$ code.
The final generated code now contains a single constructor application that is not guarded by a $\textit{reuse}$, the one for
the equation:
\begin{lstlisting}
insert E x = T R E x E
\end{lstlisting}
The generated code now has the property that if the input tree is not shared, then only a single new node is allocated.
Moreover, even if the input tree is shared we have observed a positive performance impact using \kw{reset} and \kw{reuse}.
The recursive call $(\textit{insert}\ a\ x)$ always returns a nonshared node even if $x$ is shared. Thus,
$\textit{balance}_1\ y\ b\ (\textit{insert}\ a\ x)$ always reuses at least one memory cell at runtime.

There is another way to avoid the unnecessary memory allocation that does not rely on inlining.
We can chain the $T$ constructor value from $\textit{insert}$ to $\textit{balance}_1$.
We accomplish this by rewriting $\textit{balance}_1$ and $\textit{insert}$ as follows
\begin{lstlisting}
balance_1 (T _ _ v t) (T _ (T R l x r_1) y r_2)  = T R (T B l x r_1) y (T B r_2 v t)
balance_1 (T _ _ v t) (T _ l_1 y  (T R l_2 x r)) = T R (T B l_1 y l_2) x (T B r v t)
balance_1 (T _ _ v t) (T _ l y r)               = T B (T R l y r) v t
insert (T B a y b) x = balance_1 (T B E y b) (insert a x) if x < y and a is red
\end{lstlisting}
Now, the input value $(T\ B\ a\ y\ b)$ is reused to create value $(T\ B\ E\ y\ b)$ which is passed to $\textit{balance}_1$. Note that
we have replaced $a$ with $E$ to make sure the recursive application $(\textit{insert}\ a\ x)$
may also perform destructive updates if $a$ is not shared.
This simple modification guarantees that $\textit{balance}_1$ does not allocate memory when
the input trees are not shared.

%% \subsection*{Monad transformers}
%% Monads provide a powerful way to build computations with effects~\cite{WadlerMonad},
%% and are ubiquitous in pure functional programming. Monad
%% transformers~\cite{MonadTransformers} take a monad as an argument and return a
%% new monad with new effects. Monad transformers are used to compose
%% effects encapsulated by monads in a modular way. We use monad
%% transformers extensively in the new implementation of the Lean
%% programming language.  The new Lean parser, macro expander, IR compiler, and
%% elaborator are all implemented using monad transformers.  A small
%% modification on these monad transformers allows our compiler to
%% generate much better code that avoids unnecessary memory allocations.
%% We demonstrate the main idea using the state monad, but one can apply the same idea
%% in the implementation of the state and exception monad transformers.
%% \begin{lstlisting}
%% State $\sigma$ $\alpha$ = $\sigma$ -> ($\alpha \times \sigma$)
%% pure (a : $\alpha$) : State $\sigma$ $\alpha$ = $\lambda$ s. (a, s)
%% bind (x : State $\sigma$ $\alpha$) (f : $\alpha$ -> State $\sigma$ $\beta$) : State $\sigma$ $\beta$ =
%%   $\lambda$ s. let (b, s') = x s in f b s'
%% get : State $\sigma$ $\sigma$ = $\lambda$ s. (s, s)
%% put (s' : $\sigma$) : State $\sigma$ Unit = $\lambda$ s. (unit, s')
%% \end{lstlisting}
%% The $\textit{State}$ monad allows us to chain the state
%% $\sigma$ through a series of function calls using the $\textit{bind}$
%% combinator.  The action $\textit{get}$ returns the current state, and
%% $\textit{put}\ s$ updates it. Note that the functions $\textit{pure}$,
%% $\textit{get}$ and $\textit{put}$ need to allocate a pair to store the
%% result.  This is a known problem, and Haskell uses unboxed pairs in
%% the definition of the IO monad to avoid this allocation~\cite{haskelUnboxed}. However, to
%% the best of our knowledge, unboxed pairs cannot be used in the
%% definition of polymorphic code, and monad transformer libraries cannot use them.
%% The worker/wrapper optimization~\cite{Jones96compilinghaskell} may introduce unboxed return types locally,
%% but does not do so for sum types (such as from the exception monad) or nested return
%% values (e.g.\ from a monad transformer) in its current implementation.
%% We now consider an alternative definition of the state monad where we chain not only the state, but also a pair.
%% \begin{lstlisting}
%% State $\sigma$ $\alpha$ = (Unit $\times$ $\sigma$) -> ($\alpha \times \sigma$)
%% pure a = $\lambda$ (unit, s). (a, s)
%% bind x f = $\lambda$ (unit, s). let (b, s') = x (unit, s) in f b (unit, s')
%% get = $\lambda$ (unit, s). (s, s)
%% put s' = $\lambda$ (unit, s). (unit, s')
%% \end{lstlisting}
%% Given this definition, our compiler automatically inserts $\textit{reset}$
%% and $\textit{reuse}$ instructions to reuse this pair and avoid memory
%% allocations at $\textit{pure}$, $\textit{get}$ and $\textit{put}$.
%% Our experimental results demonstrate this simple modification produces
%% a significant performance improvement.
%% We found a similar pattern in the exception monad transformer where $\textit{pure}\ a$
%% allocates a value of type $\textit{Either}$. We can change this
%% transformer to chain a value of type $\textit{Either}$, and again
%% avoid the memory allocation.


%%% Local Variables:
%%% mode: latex
%%% TeX-master: "refcount"
%%% End:
